\documentclass[UTF8,12pt,a4paper]{ctexart}

% CUMCM 2026 合规基线模板，逐条对齐官方《论文格式规范（2026年修订稿）》原文：
%   * 电子版第一页必须为摘要专用页：承诺书、编号专用页不得出现在电子版论文中（第十条）
%   * 摘要含标题和关键词，无需英文，原则上 ≤1 页；页码从摘要页起，页脚居中从 1 连续编号（第三条）
%   * 正文从第四页（纸质口径）开始，不要目录，不超过 30 页（第四条）
%   * 附录必须含支撑材料文件列表 + 全部完整可运行源程序；无程序/无支撑材料须在附录明确声明（第五条）
%   * 页边距上下左右至少 2.5cm（第一条，本模板已设 2.5cm）
%   * 任何地方不得出现参赛者身份、学校、赛区信息（第六条）
%   * 摘要/正文/附录中所有引用必须列参考文献并在正文标注（第七条）
%   * 字号、字体、行距、颜色官方不统一要求（第八条）
%   * 电子版建议 PDF，单个文件 ≤20MB，不压缩，与纸质版内容一致（第十条）
% 当届官方模板存在时以其为准复制替换本文件。
\usepackage{amsmath,amssymb}
\usepackage{booktabs}
\usepackage{geometry}
\usepackage{graphicx}
\usepackage[hidelinks]{hyperref}
\geometry{left=2.5cm,right=2.5cm,top=2.5cm,bottom=2.5cm}
\newcommand{\keywords}[1]{\par\noindent\textbf{关键词：}#1}

\title{LATEX_TEMPLATE_TITLE}
\author{}   % 匿名：不得出现姓名、学校、地区等身份信息
\date{}

\begin{document}
\maketitle

\begin{abstract}
LATEX_TEMPLATE_ABSTRACT
\keywords{LATEX_TEMPLATE_KEYWORDS}
\end{abstract}

\section{问题重述与分析}
LATEX_TEMPLATE_BODY

\section{模型假设与符号说明}
LATEX_TEMPLATE_BODY

\section{模型建立、求解与验证}
LATEX_TEMPLATE_BODY

\section{灵敏度分析与稳健性检验}
LATEX_TEMPLATE_BODY

\section{模型评价与适用边界}
LATEX_TEMPLATE_BODY

\section*{AI 工具使用声明}
% 2026 年新口径：须说明具体用途；不单独开页，紧随排版；参考文献中不再列 AI 条目。
LATEX_TEMPLATE_AI_STATEMENT

\bibliographystyle{plain}
\bibliography{references}

\appendix
\section*{支撑材料文件清单}
% 官方第五条：附录必须包含支撑材料文件列表；赛区建议置于附录第一页。
% 确实没有支撑材料时改为注明“本论文没有支撑材料”。
LATEX_TEMPLATE_SUPPORT_MANIFEST

\section{程序与补充说明}
% 官方第五条：附录必须包含建模所用全部完整、可运行的源程序代码。
% 确实没有用到程序时改为注明“本论文没有用到程序”。
LATEX_TEMPLATE_APPENDIX

\end{document}
